\section{Locality and counting}\label{sec:locality}

The basic tool of the paper is a cut lemma: if a set $J$ of time steps must produce, in the
output, bits of the input that $J$ never reads, then those bits must enter $J$ through its
\emph{entry data}, the few states and work-tape symbols that $J$ inherits from the rest of the
run. This section states the cut lemma for an arbitrary set of steps; \cref{sec:blocks} bounds the
entry data of unions of time slots, and later sections apply it to passes, trace classes and
networks.

\subsection{Entry data and cut sets}

Fix a machine $M$, an input $w$ and a finite set $J$ of time steps.

\begin{definition}[Entry data]\label{def:entry}
\formal{Tape/Defs.lean}{entryPoints, entryCells, IsLastVisit}
An \emph{entry point} of $J$ is a step $\tau\in J$ with $\tau\ne0$ and $\tau-1\notin J$. An
\emph{entry cell} of $J$ is a pair $(t,\tau)$ of a work tape $t$ and a step $\tau\in J$ such that
the cell under head $t$ at time $\tau$ was last visited by that head at a step $\sigma\notin J$
(so $\sigma<\tau$). Cells never visited before $\tau$ are blank for every input and are not entry
cells. We write
\[
  \eone(J)=\#\{\text{entry points of }J\},\qquad \etwo(J)=\#\{\text{entry cells of }J\}.
\]
\end{definition}

Step $0$ is never an entry point: the state at time $0$ is the initial state for every input.

\begin{definition}[Cut set, sources, sinks, cut pairs]\label{def:cut}
\formal{Tape/Defs.lean}{cutSet, srcRows, snkCols, cutPairs, tpos}
Let $w_0$ be an input of length $n^2b$. Relative to the run of $M$ on $w_0$:
\begin{itemize}[leftmargin=2em]
\item the \emph{cut set} of $J$ is the set of input positions $d<n^2b$ that the input head visits
at no step of $J$ but whose transposed copy, output bit $\operatorname{tpos}(d)$, is appended at a
step of $J$;
\item the \emph{source rows} $\src(J)$ are the rows $i<n$ such that the input head lies in the
row's interval $[inb,(i+1)nb)$ at some step of $J$;
\item the \emph{sink columns} $\snk(J)$ are the columns $j<n$ all of whose output bits
$[jnb,(j+1)nb)$ are appended at steps of $J$;
\item the \emph{cut pairs} are $S_J=([n]\setminus\src(J))\times\snk(J)$.
\end{itemize}
\end{definition}

Every bit of an entry $(i,j)$ with $(i,j)\in S_J$ lies in the cut set, so
$b\,\abs{S_J}\le\abs{\mathrm{cut}(J)}$ (\lean{cutPairs_card_mul_le}).

\subsection{The Locality Lemma}

\begin{lemma}[Locality]\label{lem:locality}
\formal{Tape/Statements/Core.lean}{locality, locality_oblivious, locality_function}
Let $w,w'$ be inputs whose runs have the same work-head positions at all times $\tau<T$, and let
$J\subseteq[0,T)$. Suppose the two runs agree on
\begin{enumerate}[label=(\roman*)]
\item the state at every entry point of $J$,
\item the symbol under the work head at every entry cell of $J$ (entry cells taken in the run on
$w$), and
\item the input symbol read at every step of $J$.
\end{enumerate}
Then at every step $\tau\in J$ the two runs have the same state, the same symbols under the work
heads, and take the same transition. In particular, if $M$ is oblivious on a set $W$ with halting
time $T$, then for every $J\subseteq[0,T)$ there is a function $F$ such that for every $w\in W$
the transitions taken during $J$ are $F$ applied to the states of $w$ at the entry points, the
symbols of $w$ at the entry cells, and the input symbols read by $w$ during $J$.
\end{lemma}

The run need not be live: either run may halt during $J$, since halted runs are frozen.

\begin{outline}
Induction over the steps of $J$ in increasing order. At $\tau\in J$, the state is given by~(i)
if $\tau$ is an entry point, and otherwise is produced by the common transition at $\tau-1\in J$.
The symbol under work head $t$ is given by~(ii) if $(t,\tau)$ is an entry cell; is blank in both
runs if the cell was never visited (\cref{lem:cells}); and otherwise was written at its last visit
$\sigma\in J$, where by induction both runs took the same transition at the same head position.
With~(iii), the transition function receives identical arguments. The functional form follows
because, on an oblivious set, entry points, entry cells and input positions are the same for all
inputs.
\end{outline}

The same induction gives the forms used for machines whose heads depend on the data. In those
forms the head positions are not assumed equal; they are derived.

\begin{lemma}[Locality, data-dependent heads]\label{lem:locality-core}
\formal{Tape/Statements/Core.lean}{locality_core, locality_interval}
Let $J$ be a set of steps and $A$ a set of \emph{anchors} such that every non-anchor step of $J$
has its predecessor in $J$. Suppose that two runs agree, at every anchor in $J$, on the state, all
head positions and the symbols under the work heads; that the input symbols agree wherever the
input heads agree; and that whenever the work heads agree at a non-anchor step $\tau\in J$, each
cell they read has, at some time $a\le\tau$, the same content in both runs, with the last visits
to that cell during $[a,\tau)$ in both runs lying in $J$. Then the runs agree at every step of $J$
on the state, all head positions, the symbols under the work heads and the transition taken. In
particular, two runs that agree at time $a$ on the state and all head positions, on the contents
at time $a$ of every work cell visited during $[a,b)$, and on the input bits read during $[a,b)$,
agree on all of this at every time in $[a,b]$ (transitions for $\tau<b$).
\end{lemma}

\begin{outline}
The induction of \cref{lem:locality}, carrying the head positions along: at a non-anchor step the
positions agree because the previous transitions agree. The interval form is the case
$J=[a,b)$, $A=\{a\}$.
\end{outline}

\subsection{The Cut Lemma}

\begin{lemma}[Cut Lemma]\label{lem:cut}
\formal{Tape/Statements/Core.lean}{cut_bound, cut_bound_pairs, cut_bound_bits, card_le_entryData}
Let $M$ be oblivious on $\mathrm{Mat}(n,b)$ with halting time $T$ and transpose $n\times n$
matrices with $b$-bit entries. Let $w_0\in\mathrm{Mat}(n,b)$ and $J\subseteq[0,T)$. Then
\[
  2^{\abs{\mathrm{cut}(J)}}\;\le\;(q+1)^{\eone(J)}\,(g+1)^{\etwo(J)},
  \qquad\text{hence}\qquad
  2^{b\abs{S_J}}\le(q+1)^{\eone(J)}(g+1)^{\etwo(J)}
\]
and $b\abs{S_J}\le\eone(J)\clog{(q+1)}+\etwo(J)\clog{(g+1)}$.
\end{lemma}

The lemma has a form for machines that are oblivious and correct only on a subset of the inputs.
Nothing is assumed about inputs outside the subset.

\begin{lemma}[Cut Lemma on a set]\label{lem:cut-on}
\formal{Tape/Statements/Core.lean}{cut_bound_on, cut_bound_on_coentropy}
Let $M$ be oblivious on a set $W$ with halting time $T$ and correct on $W$ at time $T$, let
$w_0\in W$ and $J\subseteq[0,T)$. Then
\[
  2^{\abs{\mathrm{cut}(J)}}\,\abs{W}\;\le\;2^{n^2b}\,(q+1)^{\eone(J)}(g+1)^{\etwo(J)} .
\]
In particular, if $W$ has co-entropy at most $m$, then
$2^{\abs{\mathrm{cut}(J)}}\le 2^m(q+1)^{\eone(J)}(g+1)^{\etwo(J)}$.
\end{lemma}

\begin{outline}
Group the inputs of $W$ by their bits outside the cut set: there are at most
$2^{n^2b-\abs{\mathrm{cut}(J)}}$ groups. Within a group the inputs agree on every input bit read
during $J$, because the cut set is read at no step of $J$ and all inputs of $W$ share the input
trajectory. By \cref{lem:locality} the transitions during $J$, and hence the output bits appended
during $J$, are a function of the entry data. These output bits contain the transposed copy of
every cut-set bit, and by correctness they equal it; so distinct inputs of a group have distinct
entry data. Hence each group has at most $(q+1)^{\eone}(g+1)^{\etwo}$ members. The formal proof
packages this injection as an abstract counting core, shared with \cref{lem:cut-trace}; the case
$W=\mathrm{Mat}(n,b)$ gives \cref{lem:cut}.
\end{outline}

\subsection{Families of step sets}

Applying the Cut Lemma to several step sets and multiplying gives the form used for partitions.

\begin{corollary}[Cut Lemma for families]\label{cor:cut-family}
\formal{Tape/Statements/Core.lean}{cut_bound_sum, cut_bound_sum_bits, cut_bound_sum_on, sum_cutPairs_ge, uncovered_cols_le, uncovered_cols_le_len}
Under the hypotheses of \cref{lem:cut}, for every finite family $(J_u)_{u\in I}$ of subsets of
$[0,T)$ (not necessarily disjoint),
\[
  b\sum_{u}\abs{S_{J_u}}\;\le\;\clog{(q+1)}\sum_u\eone(J_u)+\clog{(g+1)}\sum_u\etwo(J_u),
\]
and under the hypotheses of \cref{lem:cut-on} with co-entropy $m$ the right side of the
exponential form gains a factor $2^{\abs{I}m}$. Moreover:
\begin{enumerate}[label=(\alph*)]
\item if every $J_u$ reads at most $\rho$ rows, then
$\bigl|\bigcup_u\snk(J_u)\bigr|\,(n-\rho)\le\sum_u\abs{S_{J_u}}$;
\item if the $J_u$ are the label classes of a labelling $\lambda$ of the steps $<T$ and $b\ge1$,
then the number of columns that are not sinks of a single class is at most the number of label
changes $\#\{1\le\tau<T:\lambda(\tau)\ne\lambda(\tau-1)\}$. This holds for any run that has
appended at least $n^2b$ output bits by time $T$, with no obliviousness, correctness or halting
assumption.
\end{enumerate}
\end{corollary}

\begin{outline}
The first part multiplies the bounds of \cref{lem:cut} and \cref{lem:cut-on}. For~(a), a sink
column of $J_u$ pairs with each of the at least $n-\rho$ rows not read by $J_u$. For~(b), the
output bits of a column are appended during an interval of time; if they are not all appended in
one label class, the label changes inside that interval, and the intervals of distinct columns
are disjoint.
\end{outline}
